% 原创教学几何示意；三幅图均使用固定坐标。
\begin{tikzpicture}[figure picture,foundation picture,x=1mm,y=1mm,font=\figurefont\fontsize{8.5}{11}\selectfont]
  \foreach \shift/\title in {0/聚类,57/分类间隔,114/三元关系}{
    \begin{scope}[xshift=\shift mm]
      \node[font=\figurefont\small] at (22,44) {\title};
      \fill[FigurePaper] (0,0) rectangle (46,36);
    \end{scope}
  }
  \foreach \x/\y in {5/6,8/17,16/7,20/14}{
    \draw[FigureBlue!45] (\x,\y)--(12,11);
    \fill[FigureBlue] (\x,\y) circle (1);
  }
  \foreach \x/\y in {26/22,31/31,40/25,38/18}{
    \draw[FigureRed!50] (\x,\y)--(34,24);
    \fill[FigureRed] (\x-0.9,\y-0.9) rectangle (\x+0.9,\y+0.9);
  }
  \foreach \x/\y in {12/11,34/24}{
    \draw[FigureStroke,line width=1.1pt] (\x-2,\y)--(\x+2,\y) (\x,\y-2)--(\x,\y+2);
  }
  \begin{scope}[xshift=57mm]
    \fill[FigurePanel] (18,0) rectangle (28,36);
    \draw[FigureStroke,line width=1.1pt] (23,0)--(23,36);
    \draw[FigureStroke,dashed] (18,0)--(18,36) (28,0)--(28,36);
    \foreach \x/\y in {6/8,10/22,15/30,18/15}{
      \fill[FigureBlue] (\x,\y) circle (1);
    }
    \foreach \x/\y in {28/9,32/25,37/17,40/31}{
      \fill[FigureRed] (\x-0.9,\y-0.9) rectangle (\x+0.9,\y+0.9);
    }
    \draw[<->,FigureMuted] (18,39)--(28,39);
  \end{scope}
  \begin{scope}[xshift=114mm]
    \draw[FigureRed,line width=1.1pt] (9,10)--(18,24);
    \draw[LancetGreen,line width=1.1pt,dashed] (9,10)--(37,17);
    \fill[FigureBlue] (9,10) circle (1.4);
    \fill[FigureRed] (18,24) circle (1.4);
    \fill[LancetGreen] (36,16) rectangle (38,18);
    \node[below] at (9,8) {$\bm a$};
    \node[above] at (18,26) {$\bm p$};
    \node[above] at (37,19) {$\bm n$};
  \end{scope}
\end{tikzpicture}
